This chapter interprets the empirical results beyond metric reporting, connects them back to the research questions, and discusses the practical and methodological implications of treating predictive multiplicity as a spatially structured phenomenon.

\subsection{From Global Multiplicity to Local Risk}

The central empirical message of this thesis is that predictive multiplicity is \emph{localized}: it is not spread uniformly across the population, but concentrates in identifiable regions of feature space. Global disagreement metrics such as ambiguity, discrepancy, or disagreement rate remain useful because they quantify \emph{how much} disagreement exists overall. However, by construction, they do not reveal \emph{where} disagreement occurs or \emph{which} individuals are most exposed to it. By combining observation-wise predictive variance with Moran's $I$ and LISA, this thesis shows that disagreement is often organized into contiguous HH (High--High) regions---multiplicity hotspots---rather than being randomly scattered.

This changes the practical interpretation of predictive reliability. Under a single deployed model, two individuals may receive similarly confident predicted probabilities, yet their predictions need not be equally reliable. One individual may lie in a stable region where near-optimal models agree closely, while the other may lie in a hotspot where equally accurate models disagree substantially. In that case, the second prediction is operationally less reliable even when its reported score appears equally confident. The key point is therefore not only \emph{how uncertain} a prediction is, but also \emph{how contingent it is on arbitrary model choice within the Rashomon set}.

\paragraph{Hotspots form regions, not isolated points.}
The connected-component analysis (Section~\ref{sec:results-spatial}) shows that HH observations are not merely isolated outliers. On COMPAS, hotspot points typically form a small number of coherent connected regions, with $3.7 \pm 2.0$ components per run and a mean largest component of $36.9 \pm 22.7$ points. Predictive multiplicity is therefore regional rather than point-wise random. This matters in practice because regional structure is more actionable than isolated unstable predictions: it enables subgroup-level auditing, rule extraction, and targeted monitoring.

\paragraph{Localized multiplicity is different from standard uncertainty reporting.}
It is also important to distinguish the present framework from more standard uncertainty analyses. Ensemble variance, Bayesian posterior uncertainty, or repeated-training variability usually quantify uncertainty around one modeling pipeline or one stochastic training process. Rashomon multiplicity asks a different question: \emph{how much do equally acceptable models disagree, and where does that disagreement concentrate?} A point can therefore be confidently predicted by each model individually and still be multiplicity-sensitive if acceptable models disagree systematically with one another. The hotspot analysis adds a local diagnostic layer that conventional uncertainty summaries do not provide.

\subsection{Research Questions Revisited}

\paragraph{RQ1--RQ2 (prediction-level multiplicity and spatial structure).}
Across datasets, observation-wise predictive variance is nontrivial and often spatially clustered. The strongest evidence appears on COMPAS, which shows the largest Moran's $I$ and the clearest hotspot structure, while German Credit exhibits substantial multiplicity with weaker but still meaningful spatial organization. Adult shows lower spatial concentration than COMPAS but still a positive and statistically supported clustering signal. The permutation-based null experiments show that this structure is not well explained by random reassignment of variance values to observations. Together, these results answer the first two research questions positively: predictive multiplicity is not only present at the observation level, but is often spatially organized in feature space.

\paragraph{RQ3 (stability of spatial patterns).}
The stability results reveal an important distinction between point-level and region-level reproducibility. Exact HH point sets overlap only weakly across seeds (mean pairwise Jaccard $\sim$0.025 on COMPAS), and only a small fraction of observations are repeatedly labeled as HH. At the same time, hotspot \emph{regions} are markedly more stable than raw HH masks. When connected components are compared across runs, the same broader areas of feature space recur much more consistently (mean best-match Jaccard across components $\sim$0.35). Thus, the precise hotspot boundary is unstable, but the broader geographic location of instability is comparatively persistent. This supports a region-based interpretation of multiplicity: the exact membership of a hotspot may shift from run to run, but the same neighborhoods of feature space repeatedly emerge as locally fragile.

\paragraph{RQ4 (drivers of multiplicity).}
The variance decomposition shows that model family is a major source of disagreement, accounting for a substantial fraction of predictive multiplicity globally and an even larger share within HH hotspots. This indicates that the choice between model classes is especially consequential precisely in those regions where the system is least stable. Within families, multiplicity is further shaped by a relatively small number of hyperparameters: regularization strength for logistic regression, neighborhood size for $k$NN, depth and learning-rate controls for tree-based models, and optimization and architecture choices for MLPs. Predictive multiplicity is therefore not a diffuse byproduct of random tuning; it is structured by concrete modeling decisions.

\paragraph{RQ5 (robustness and sensitivity).}
The main spatial findings persist across a broad set of analytical choices, including different Rashomon set sizes $K$, different neighborhood sizes $k$, post-hoc calibration, and alternative graph constructions such as PCA-based or cosine-distance $k$NN graphs. These checks matter because hotspot analyses can otherwise be dismissed as artifacts of a particular graph or threshold choice. Here, the qualitative conclusions remain stable even when effect sizes shift. This supports the interpretation that the detected multiplicity structure reflects properties of the data-model interaction rather than a fragile tuning decision.

\paragraph{RQ6--RQ7 (synthetic validation and interpretability).}
The synthetic experiments show that the pipeline is capable of recovering known ambiguous regions with high precision and recall, which provides an important validation of the overall approach. The rule extraction results further show that hotspot regions can be summarized by compact feature conditions, especially on COMPAS. This makes the framework practically useful: it does not merely state that instability exists, but also identifies which subpopulations are most affected and how they can be described in human-interpretable terms.

\subsection{Practical Implications}

The results suggest a deployment workflow that differs from standard ``single best model'' practice:
\begin{enumerate}
    \item train and retain a near-optimal set of models rather than only one selected winner;
    \item compute pointwise multiplicity metrics and hotspot maps on held-out data;
    \item treat predictions in hotspot regions differently, for example by attaching caution labels, triggering human review, or using abstention policies in high-stakes settings.
\end{enumerate}

The practical contribution is therefore not merely another descriptive metric, but a way to separate \emph{stable predictions} from \emph{model-contingent predictions}. In high-stakes applications, this distinction is operationally relevant. A system should not treat a prediction in a stable region and a prediction in a hotspot as equally trustworthy simply because both come with similar point estimates.

A useful way to interpret the proposed workflow is as a \emph{reliability triage} mechanism. Predictions in stable regions may proceed through the standard pipeline. Predictions in localized hotspots should be interpreted more cautiously because the selected model is standing in for many equally accurate alternatives that would not necessarily agree. The precise intervention will depend on the application domain, but the general principle is clear: multiplicity hotspots identify cases where model choice itself becomes decision-relevant.

The subgroup exposure analysis also suggests that multiplicity can be unevenly distributed across demographic groups. Even though the estimated disparities are not consistently significant across all seeds, they are substantial in some runs and therefore important from a governance perspective. This motivates routine multiplicity audits as a complement to standard fairness monitoring. A model can satisfy conventional performance-based fairness checks while still exposing some groups more strongly to model arbitrariness.

\subsection{Methodological Reflection}

A central conceptual takeaway is that predictive multiplicity and decision-boundary uncertainty are related, but not identical. On Adult and German Credit, high-variance observations tend to lie closer to the decision boundary, which suggests that multiplicity partly reflects ordinary classification ambiguity. On COMPAS, by contrast, the relationship between variance and boundary proximity is weak. There, multiplicity appears to arise less from intrinsic closeness to the boundary and more from structural differences between equally accurate model families and hyperparameter settings.

This distinction matters because it clarifies what kind of phenomenon the proposed pipeline is detecting. It is useful to separate three perspectives:
\begin{itemize}
    \item \textbf{aleatoric or data ambiguity}, arising from intrinsic class overlap;
    \item \textbf{epistemic or model uncertainty}, arising from limited data or imperfect estimation;
    \item \textbf{Rashomon multiplicity}, arising because many behaviorally different models remain equally acceptable under the chosen performance criterion.
\end{itemize}
The present framework primarily targets the third category. It is therefore best understood not as a replacement for uncertainty quantification, but as a complementary diagnostic that asks whether predictions depend strongly on which acceptable model happens to be selected.

A second reflection concerns when the method is most informative. If Moran's $I$ is near zero and hotspot regions do not emerge, then multiplicity may still be present, but it is diffuse rather than localized. In that case, global disagreement summaries may be more informative than spatial subgroup descriptions. Conversely, if an extremely large share of the dataset becomes hotspot-like, this would indicate broad instability rather than a small set of localized risk regions. The method is therefore especially valuable in the intermediate regime identified in this thesis: disagreement exists, but it is concentrated enough to support localization, explanation, and targeted intervention.

\subsection{Limitations and Threats to Validity}

Despite the strong empirical evidence, several limitations remain.

\begin{itemize}
    \item \textbf{Finite approximation of the Rashomon set.} The top-$K$ candidate approach approximates a potentially much larger and more continuous set of acceptable models. The observed multiplicity structure therefore depends on the candidate families and hyperparameter spaces explored in the experiments.
    
    \item \textbf{Dependence on the graph construction.} Although robustness checks across alternative graph constructions are encouraging, the spatial analysis still depends on how neighborhood structure is defined in transformed feature space. Different representations or manifold-aware graph constructions could reveal somewhat different local patterns.
    
    \item \textbf{Boundary instability of local hotspot labels.} Exact HH membership is unstable across seeds, even when broader hotspot regions recur. This means the method is better suited to identifying unstable \emph{areas} than to declaring a single observation intrinsically and permanently hotspot-like.
    
    \item \textbf{Small subgroup uncertainty.} Some demographic categories contain few observations, which leads to large variability in subgroup-specific exposure rates and limits inferential strength.
    
    \item \textbf{Descriptive fairness analysis.} The subgroup exposure results are observational and do not identify causal mechanisms behind disparity. They should therefore be interpreted as auditing signals, not causal claims.
    
    \item \textbf{External validity.} The experiments cover three benchmark datasets and synthetic validation tasks. The practical usefulness of multiplicity hotspot analysis should be further tested on larger, domain-specific, and operational datasets before drawing broad deployment conclusions.
\end{itemize}

\subsection{Position Relative to Prior Work}

Relative to prior work on predictive multiplicity, the main contribution of this thesis is to add a spatially explicit diagnostic layer. Existing approaches largely quantify disagreement globally or per observation, but they typically stop short of asking whether disagreement forms coherent local regions in feature space. The present results show that adding this localization step changes the practical interpretation of multiplicity: the issue is not only that models disagree, but that they disagree in systematic, recurring, and interpretable subregions of the population.

This perspective also complements interpretability research. Much of the interpretability literature explains why a \emph{single} selected model behaves as it does. Here the focus shifts to a different question: \emph{where does model choice itself matter?} In governance settings, this can be the more relevant question. When many models are equally defensible, explaining one chosen model is not sufficient if another equally good model would treat the same individuals differently.

\subsection{Future Directions}

Several extensions follow naturally from the present results:
\begin{itemize}
    \item continuous or optimization-based approximations of the Rashomon set rather than finite top-$K$ candidate pools;
    \item multiplicity-aware selective prediction, abstention, or escalation policies for high-stakes use;
    \item temporal tracking of hotspot drift under distribution shift or concept drift;
    \item stronger inferential fairness analyses using hierarchical, longitudinal, or causal designs;
    \item extension of the framework to multiclass classification and regression.
\end{itemize}

Overall, the evidence supports the central thesis claim: predictive multiplicity is not merely random training noise or a global nuisance summary, but a structured property of model classes interacting with data geometry. In high-stakes settings, it should therefore not only be measured globally, but also localized, interpreted, and reported.